\uuid{4fBy}
\exo7id{3649}
\titre{exo7 3649}
\auteur{quercia}
\organisation{exo7}
\datecreate{2010-03-10}
\isIndication{false}
\isCorrection{false}
\chapitre{Endomorphisme particulier}
\sousChapitre{Autre}
\cpge{mpsi-mp2i.dimension-finie.a, pcsi.dimension-finie.a, ptsi.dimension-finie.a, mpsi-mp2i.applications-lineaires.c, bcpst1.espace-vectoriel-kn.b, bcpst2.espaces-vectoriels.b}
\module{Algèbre}
\niveau{L2}
\difficulte{}

\contenu{
\question{
Soit $E$ un $ K$-ev de dimension finie $n$,
${\cal F} = ({\vec e}_1,\dots, {\vec e}_{n-1})$ une famille libre de $E$
et $f \in E^*\setminus\{0\}$.
Montrer qu'on peut compléter $\cal F$ en une base
${\cal B} = ({\vec e}_1,\dots,{\vec e}_n)$
telle que $f = e_n^*$ si et seulement si
$f({\vec e}_1) = \dots = f({\vec e}_{n-1}) = 0$.

Y~a-t-il unicité de ${\vec e}_n$ ?
}
}
